\section{Discussion}
\label{sec:discussion}

\paragraph{When graph methods win.}
Graph-guided retrieval shows its strongest advantage on repositories with well-defined structural dependencies---particularly Flask, where import chains and class hierarchies directly connect issue-relevant files. The graph's typed edges enable multi-hop navigation that reaches structurally related files even when their text content has low embedding similarity to the issue description. The setup cost advantage is also most pronounced in these cases: the graph index embeds only structural metadata (names, signatures, docstrings), requiring roughly one-third the embedding tokens of a full code chunk index.

\paragraph{When RAG methods win.}
On Requests and Pytest, RAG methods outperform graph methods on retrieval quality. Both repositories contain issues whose resolution involves files that are semantically related to the issue text but not strongly connected via import or call edges in the graph. This suggests that graph navigation is less effective when the structural signal is weak or when the issue description uses natural language that maps better to code semantics than to structural topology. The progressive RAG agent's ability to iteratively refine its search also helps in cases where the initial retrieval set misses relevant files.

\paragraph{Cost-quality frontier.}
The most striking finding is the separation between retrieval quality and total cost. GM-progressive and RAG-progressive achieve similar macro F1 across the three completed SWE-bench repositories (0.684 vs.\ 0.680), but GM-progressive costs 0.42$\times$ the total tokens. This cost advantage arises almost entirely from lower setup cost (0.31$\times$), partially offset by higher runtime LLM usage (2.4$\times$). As the number of tasks per snapshot increases, the amortized cost gap widens further in favor of graph methods.

\paragraph{Deterministic vs.\ LLM-guided retrieval.}
The deterministic graph scorer (\texttt{gm\_deterministic}) provides a zero-LLM-runtime baseline that is fully reproducible and auditable. Preliminary results suggest it is competitive with LLM-guided retrieval on structurally regular repositories, though LLM-guided agents retain an advantage when issues require adaptive exploration that the fixed BFS expansion cannot anticipate. A full comparison across the evaluation matrix is pending.

\paragraph{Amortization structure.}
The dual-track evaluation reveals that amortization opportunity is highly dependent on experimental design. Under strict commit fidelity, most SWE-bench issue sets have \texttt{commit\_repeat\_ratio} near zero for small $n$, eliminating any index reuse. Under same-snapshot settings, the graph setup cost drops by an order of magnitude ($\sim$12$\times$ for \texttt{psf/requests}). This structural observation is important for practical deployment: graph-based retrieval is most cost-effective in scenarios where a repository snapshot serves multiple queries, such as issue triage, code review, or batch analysis.

\paragraph{Limitations of current evidence.}
The results presented here are based on a limited number of repositories and issues, with the full evaluation matrix still partially incomplete. Key limitations include: (1) small sample sizes ($n=10$ per repo) that limit statistical power; (2) evaluation on a single LLM backend (Gemini), leaving open the question of whether results transfer to other models; (3) retrieval-only evaluation without downstream patch generation, meaning we cannot yet confirm that retrieval quality improvements translate to end-to-end issue resolution success; and (4) Python-only graph construction, which constrains generalization to multi-language repositories.
